\section{Introduction}
\label{sec:introduction}

In order to prevent duplicating work and wasting resources, researchers must verify their hypotheses against existing literature. As the volume of scientific publications has grown, this verification has become substantially harder. As of December 2025, arXiv lists 2,911,812 submissions and receives more than 26,000 new paper submissions per month \citep{arxivStats}. Since its launch in 1991, arXiv has expanded to a scale at which researchers often struggle to read and assess all relevant work before starting a new study. The rapid growth of academic publishing has produced a clear need for tools that can assess the originality of research consistently and reliably. Manual literature reviews are time consuming and frequently left incomplete, and they rarely identify precise similarities between new concepts and earlier work. They also tend to overlook prior art that lives outside academic preprints, such as patent records and public code repositories, which often document the same ideas as research papers but in different forms.

The rapid development of Large Language Models (LLMs) has made automated validation of research ideas tractable, especially for natural language understanding and information retrieval tasks. Recent work shows that LLMs are already used effectively in literature review \citep{eger2025}. However, single LLMs do not have the analytical depth required to assess the novelty of a research idea reliably. Single-prompt models perform basic comparisons and remain vulnerable to hallucination, in contrast to autonomous multi-agent systems that decompose the task across specialised roles \citep{zheng2025}. Standard retrieval systems return related papers but do not produce a structured judgement on novelty. To address this gap, our approach builds an autonomous environment in which multiple agents perform retrieval, paper-by-paper scrutiny, and global synthesis. A growing body of recent work has begun to frame literature-grounded research idea novelty as a distinct task, moving beyond simple semantic similarity toward retrieval-backed and more explainable assessment workflows \citep{shahid2025,lin2025schnovel,scholareval,rinobench,opennovelty,dasilva2025graphmind}, and beyond pure novelty checking toward broader research-assistant systems for the scientific lifecycle \citep{eger2025survey,wang2024scimon,si2024llmideas,lu2024aiscientist,baek2024researchagent}.

This study presents Hypothetica \citep{hypotheticaGitHub}, an AI-driven research validation system that evaluates the originality of a proposed research concept against multiple sources of prior art. Rather than restricting itself to a single repository, Hypothetica retrieves evidence through a pluggable adapter layer that unifies scholarly works from OpenAlex \citep{openalexDocs,priem2022openalex} and arXiv \citep{arxivAPI}, patent records from Google Patents, and open-source repositories from GitHub under a single originality-scoring pipeline. The system generates originality reports using a multi-agent architecture driven by Google Gemini \citep{geminiAPI} and a Retrieval Augmented Generation (RAG) architecture \citep{lewis2020rag,gao2024ragsurvey}, so researchers can make informed decisions about their concept on the basis of evidence drawn from across the academic, industrial, and software ecosystems.

The analysis pipeline is divided into two layers and assesses each research concept along four dimensions: novelty of the technical problem, methodological innovation, overlap in the application domain, and similarity of the claimed contribution. It integrates RAG with originality scoring, semantic chunking, and dense vector retrieval based on E5 embeddings \citep{wang2022e5} and Chroma \citep{chroma}, followed by cross-encoder reranking \citep{nogueira2019monobert} before any candidate reaches the scoring layer. Before retrieval, a conversational interview agent elicits missing details about the idea over a short multi-round dialogue, and a reality-check agent runs in parallel to flag whether the proposed concept already exists as a known product or system, independently of what the corpus returns. For source types where document-level granularity is more appropriate than passage-level retrieval, such as GitHub README files, the same scoring layer consumes the document as a whole rather than as embedded chunks. In addition to a numerical originality score, the system features a report generation component that links specific sentences in the user's research idea to the relevant sections of the retrieved sources considered in the report.

The primary contributions of this study are as follows. First, we introduce Hypothetica, an automated system designed to transform manual originality assessment into an LLM-based computational pipeline. Second, we propose a specialised multi-agent architecture that improves assessment quality compared to single-prompt baselines by assigning dedicated agents to specific roles, including a conversational interview agent for clarifying the idea, a reality-check agent for surfacing already-existing products, query-variance and relevance-filtering agents for retrieval, and per-paper Layer~1 and global Layer~2 agents for scoring and aggregation. Third, we design a pluggable evidence-adapter layer that lets the same scoring pipeline reason over heterogeneous prior art (academic works from OpenAlex and arXiv, patent records from Google Patents, and open-source repositories from GitHub) without changes to the scoring agents themselves. Fourth, instead of merely producing an overall evaluation score, our system shows exactly which parts of the user's idea match existing work by linking sentences to specific passages in the retrieved sources.
